% ROZDZIAL 12 — Model obiektowy C++ i polimorfizm
% Fragment do zbiorczego skryptu; wczytac przez \input{...} PO rozdziale 11.
% Standard podstawowy: C++20. Wymagane pakiety: listings, amsmath, amssymb,
% booktabs, tabularx, array, hyperref, xcolor (jak we wczesniejszych rozdzialach).
% Ten plik celowo nie zawiera \documentclass ani \begin{document}.

\chapter{Model obiektowy C++ i polimorfizm}
\label{chap:cpp-object-model}

\section{Po co poznawać model obiektowy?}

Klasy w C++ są jednocześnie narzędziem modelowania danych, organizowania
interfejsów i budowania różnych rodzajów polimorfizmu. Dwie implementacje
tego samego interfejsu mogą jednak mieć bardzo różne własności:
\begin{itemize}
\item mogą rozstrzygać wywołanie już w trakcie kompilacji albo dopiero w
  czasie wykonania;
\item mogą wymagać alokacji dynamicznej lub działać bez niej;
\item mogą przechowywać obiekty bezpośrednio, przez referencję albo przez
  wskaźnik zarządzający czasem życia;
\item mogą zachowywać tożsamość dynamicznego typu obiektu lub ją utracić
  wskutek \emph{object slicing}.
\end{itemize}

W tym rozdziale rozdzielamy \emph{gwarancje standardu C++} od typowych
technik używanych przez kompilatory, takich jak tablice funkcji wirtualnych
(\emph{vtable}). Pozwala to czytać kod zarówno od strony języka, jak i
możliwych kosztów wykonania.

\section{Obiekty, podobiekty i tożsamość typu}
\label{sec:cpp12-subobjects}

\emph{Obiekt} zajmuje pewien obszar pamięci, ma określony typ i czas życia.
Obiekt klasy może zawierać \emph{podobiekty składowych} oraz
\emph{podobiekty klas bazowych}. Dziedziczenie nie oznacza zatem
,,zamiany'' klasy pochodnej na bazową: obiekt pochodny zawiera w sobie
odpowiedni podobiekt bazowy.

\begin{lstlisting}[caption={Klasa pochodna zawiera podobiekt bazowy (C++11); kompletny program.}]
#include <cassert>
#include <string>

struct Entity {
    int id = 0;
};

struct Employee : Entity {
    std::string name;
};

int main() {
    Employee e;
    e.id = 7;
    e.name = "Ada";
    Entity& base = e;  // odniesienie do podobiektu Entity
    assert(&base == static_cast<Entity*>(&e));
    assert(base.id == 7);
}
\end{lstlisting}

\paragraph{Typ statyczny a typ dynamiczny.}
\emph{Typ statyczny} wyrażenia wynika z reguł języka i jest znany
kompilatorowi. Gdy wyrażenie wskazuje na obiekt klasy polimorficznej
(czyli klasy posiadającej lub dziedziczącej funkcję wirtualną),
\emph{typ dynamiczny} obiektu może być klasą pochodną.

\begin{lstlisting}[caption={Statyczny i dynamiczny typ obiektu; fragment.}]
struct Animal { virtual ~Animal() = default; };
struct Dog : Animal {};

Dog dog;
Animal& ref = dog;  // statyczny typ: Animal&; obiekt: Dog
Animal* ptr = &dog; // statyczny typ: Animal*; obiekt: Dog
\end{lstlisting}

\noindent W C++ nie ma wymogu, aby podobiekt klasy bazowej zajmował
\emph{pierwsze} bajty obiektu pochodnego. Nie należy również zakładać,
że wskaźnik do całego obiektu i wskaźnik do dowolnej jego bazy mają
tę samą reprezentację liczbową. Odpowiednie konwersje wykonuje język.

\subsection{Składnia dziedziczenia i kontrola dostępu}

\begin{lstlisting}[caption={Dziedziczenie publiczne, chronione i prywatne; fragment.}]
class Base {
public:
    int a = 1;
protected:
    int b = 2;
private:
    int c = 3;
};

class PublicDerived    : public Base {};
class ProtectedDerived : protected Base {};
class PrivateDerived   : private Base {};
\end{lstlisting}

Przy dziedziczeniu publicznym publiczne elementy interfejsu bazowego
pozostają publiczne, a chronione pozostają chronione. Przy dziedziczeniu
chronionym oba rodzaje stają się chronione, natomiast przy prywatnym
oba stają się prywatne dla użytkowników klasy pochodnej.
Prywatne elementy bazy pozostają niedostępne \emph{bezpośrednio}
wewnątrz klasy pochodnej niezależnie od rodzaju dziedziczenia.

Sama dostępność elementu nie zmienia jego faktycznego istnienia
w podobiekcie bazowym. Publiczne dziedziczenie najczęściej wyraża
relację \emph{is-a}: obiekt pochodny może być używany tam, gdzie
oczekiwany jest interfejs bazowy.

\section{Przeciążanie, przesłanianie i funkcje wirtualne}
\label{sec:cpp12-virtual}

\emph{Przeciążanie} (\emph{overloading}) polega na deklarowaniu kilku
funkcji o tej samej nazwie, różniących się listą parametrów.
\emph{Przesłanianie nazwy} (\emph{name hiding}) zachodzi, gdy deklaracja
w klasie pochodnej zasłania nazwę z klasy bazowej.
\emph{Nadpisywanie} (\emph{overriding}) oznacza dostarczenie implementacji
funkcji wirtualnej z bazy przez klasę pochodną.

\begin{lstlisting}[caption={Dynamiczne wiązanie wywołania przez virtual (C++11); kompletny program.}]
#include <iostream>

struct Shape {
    virtual ~Shape() = default;
    virtual double area() const = 0;
    virtual void print() const {
        std::cout << "Shape\n";
    }
};

struct Square final : Shape {
    explicit Square(double side) : side_(side) {}
    double area() const override { return side_ * side_; }
    void print() const override {
        std::cout << "Square: " << area() << '\n';
    }
private:
    double side_;
};

int main() {
    Square square{3.0};
    Shape& shape = square;
    shape.print(); // Square: 9
    std::cout << shape.area() << '\n';
}
\end{lstlisting}

Klasa \texttt{Shape} jest \emph{abstrakcyjna}, ponieważ jej funkcja
\texttt{area()} jest czysto wirtualna, co oznaczamy \texttt{= 0}.
Nie można skonstruować obiektu \texttt{Shape}, lecz można utworzyć
referencję lub wskaźnik do jej podobiektu w obiekcie konkretnym.

\paragraph{Reguła wyboru funkcji.}
Dla wywołania \texttt{base.f()} najpierw rozstrzygane są dostępność
nazwy, przeciążenia i poprawność wywołania według \emph{typu statycznego}.
Jeśli wybrana funkcja jest wirtualna, a wywołanie nie wyłącza mechanizmu
wirtualnego przez kwalifikację nazwy, uruchamiany jest odpowiedni
\emph{final overrider} dla dynamicznego typu obiektu.

\subsection{Dlaczego stosować \texttt{override} i \texttt{final}?}

\texttt{override} oznacza, że dana metoda ma nadpisywać metodę wirtualną.
Jeśli sygnatura nie pasuje, kompilator zgłosi błąd. \texttt{final}
uniemożliwia dalsze nadpisywanie metody albo dziedziczenie po klasie.

\begin{lstlisting}[caption={override wykrywa niezgodność sygnatur; fragment.}]
struct Base {
    virtual void draw() const {}
    virtual ~Base() = default;
};

struct Good : Base {
    void draw() const override {} // poprawnie
};

// BLEDNE -- celowo niekompilowalne:
// struct Bad : Base {
//     void draw() override {} // brak const
// };
\end{lstlisting}

\noindent Przy nadpisywaniu znaczenie mają między innymi kwalifikatory
\texttt{const}/\texttt{volatile} metody i kwalifikator referencyjny
\texttt{\&}/\texttt{\&\&}. Specyfikacja wyjątków też podlega regułom
zgodności: nadpisanie metody \texttt{noexcept} metodą potencjalnie
zgłaszającą wyjątek jest niedozwolone. Zwykłego typu wyniku nie można
zmienić, z wyjątkiem określonych \emph{kowariantnych} typów zwracanych.

\begin{lstlisting}[caption={noexcept i kwalifikatory referencyjne w nadpisywaniu; fragment.}]
struct Service {
    virtual int read() const & noexcept { return 1; }
    virtual ~Service() = default;
};

struct FastService : Service {
    int read() const & noexcept override { return 2; }
};
\end{lstlisting}

\subsection{Przesłanianie przeciążeń i \texttt{using}}

\begin{lstlisting}[caption={Odsłanianie przeciążeń klasy bazowej; kompletny program.}]
#include <cassert>

struct Printer {
    int print(int) const { return 1; }
    int print(double) const { return 2; }
};

struct FancyPrinter : Printer {
    using Printer::print;
    int print(const char*) const { return 3; }
};

int main() {
    FancyPrinter p;
    assert(p.print(10) == 1);
    assert(p.print(2.5) == 2);
    assert(p.print("text") == 3);
}
\end{lstlisting}

Bez deklaracji \texttt{using Printer::print;} nazwa \texttt{print}
z klasy pochodnej przesłaniałaby zestaw przeciążeń o tej nazwie
z klasy bazowej. Mechanizm przesłaniania występuje również dla
funkcji \emph{niewirtualnych}.

\subsection{Argumenty domyślne i wiązanie statyczne}

Wartości domyślnych argumentów nie wybiera się polimorficznie:
są związane ze statycznym typem wyrażenia, mimo że sama implementacja
funkcji może być wybrana dynamicznie.

\begin{lstlisting}[caption={Domyślny argument pochodzi z deklaracji bazy (C++11); kompletny program.}]
#include <cassert>

struct Base {
    virtual int calculate(int x = 10) const { return x; }
    virtual ~Base() = default;
};

struct Derived : Base {
    int calculate(int x = 20) const override { return 2 * x; }
};

int main() {
    Derived object;
    Base& base = object;
    assert(base.calculate() == 20); // 2 * 10
    assert(object.calculate() == 40); // 2 * 20
}
\end{lstlisting}

Z tego powodu w praktycznych interfejsach polimorficznych warto unikać
różnych domyślnych argumentów dla nadpisywanych funkcji.

\section{Destruktor wirtualny i reguły czasu życia}
\label{sec:cpp12-lifetime}

Jeżeli obiekt klasy pochodnej ma być niszczony przez wskaźnik do
klasy bazowej, destruktor bazy powinien być \texttt{public virtual}
(alternatywnie projekt może zabraniać usuwania przez bazę,
np. za pomocą \texttt{protected} destruktora niewirtualnego).

\begin{lstlisting}[caption={Bezpieczne usuwanie obiektu przez wskaźnik bazowy (C++14); kompletny program.}]
#include <iostream>
#include <memory>

struct Resource {
    virtual ~Resource() { std::cout << "~Resource\n"; }
};

struct FileResource : Resource {
    ~FileResource() override {
        std::cout << "~FileResource\n";
    }
};

int main() {
    std::unique_ptr<Resource> object =
        std::make_unique<FileResource>();
    // Najpierw ~FileResource, potem ~Resource.
}
\end{lstlisting}

Jeżeli bazowy destruktor nie jest wirtualny, standardowe
\texttt{delete base\_ptr;} przy dynamicznym typie różnym od statycznego
może prowadzić do niezdefiniowanego zachowania; istnieją szczególne
reguły związane z tzw. \emph{destroying delete}, ale nie należy
traktować ich jako zwykłego substytutu destruktora wirtualnego.

\subsection{Kolejność konstrukcji i destrukcji}

Dla obiektu klasy pochodnej kolejność konstrukcji jest następująca:
\begin{enumerate}
\item wirtualne podobiekty bazowe (przez najbardziej pochodny obiekt);
\item niewirtualne podobiekty bazowe w kolejności ich deklaracji;
\item składowe klasy w kolejności \emph{deklaracji pól}, niezależnie od
      kolejności w liście inicjalizacyjnej;
\item ciało konstruktora klasy.
\end{enumerate}
Destrukcja przebiega w odpowiedniej kolejności odwrotnej.
W prostym pojedynczym dziedziczeniu można to zilustrować następująco.

\begin{lstlisting}[caption={Kolejność konstruowania podobiektów (C++11); kompletny program.}]
#include <iostream>

struct A {
    A()  { std::cout << "A+ "; }
    ~A() { std::cout << "A- "; }
};
struct B {
    B()  { std::cout << "B+ "; }
    ~B() { std::cout << "B- "; }
};
struct C : A {
    B member;
    C()  { std::cout << "C+ "; }
    ~C() { std::cout << "C- "; }
};

int main() {
    C object;
    // A+ B+ C+ C- B- A-
}
\end{lstlisting}

\subsection{Funkcje wirtualne w konstruktorach i destruktorach}

Podczas wykonywania konstruktora \texttt{Base} wywołanie wirtualne
nie trafia do implementacji z \texttt{Derived}, której część nie
została jeszcze skonstruowana. Analogicznie podczas wykonywania
destruktora \texttt{Base} część pochodna została już zniszczona.

\begin{lstlisting}[caption={Dynamiczne wiązanie podczas konstrukcji (C++11); kompletny program.}]
#include <cassert>

struct Base {
    Base() { assert(code() == 1); }
    virtual ~Base() = default;
    virtual int code() const { return 1; }
};
struct Derived : Base {
    int code() const override { return 2; }
};
int main() {
    Derived object;
    Base& b = object;
    assert(b.code() == 2);
}
\end{lstlisting}

\noindent Nie należy wywoływać funkcji \emph{czysto wirtualnej}
przez mechanizm wirtualny z konstruktora albo destruktora klasy,
na której poziomie funkcja jest czysto wirtualna: takie wywołanie
jest niezdefiniowanym zachowaniem. Preferuj dwufazową konstrukcję
kontrolowaną przez fabrykę tylko wtedy, gdy naprawdę jest potrzebna,
a w typowych sytuacjach inicjalizację kompletnych podobiektów
przez konstruktory i RAII.

\section{Object slicing: utrata części pochodnej}
\label{sec:cpp12-slicing}

\emph{Object slicing} nie jest dynamicznym rzutowaniem. Gdy obiekt
pochodny kopiujemy \emph{przez wartość} do obiektu bazy,
powstaje osobny obiekt bazy, zawierający tylko część bazową.

\begin{lstlisting}[caption={Kopiowanie przez wartość powoduje slicing (C++11); kompletny program.}]
#include <cassert>

struct Base {
    virtual ~Base() = default;
    virtual int value() const { return 1; }
};
struct Derived : Base {
    int value() const override { return 2; }
};

int by_value(Base object) { return object.value(); }
int by_reference(const Base& object) {
    return object.value();
}

int main() {
    Derived d;
    assert(by_value(d) == 1);      // nowy obiekt Base
    assert(by_reference(d) == 2);  // oryginalny Derived
}
\end{lstlisting}

Zwykle przechowujemy heterogeniczne obiekty polimorficzne
jako \texttt{std::unique\_ptr<Base>} (własność wyłączna),
\texttt{std::shared\_ptr<Base>} (współdzielona własność) albo
nieposiadające obiektu \texttt{Base\&}/\texttt{Base*} (widoki).
Bezpośredni \texttt{std::vector<Base>} przechowuje wyłącznie obiekty
\texttt{Base} i nie zachowuje typów pochodnych przy kopiowaniu.

\subsection{Bezpieczny kontener polimorficzny}

\begin{lstlisting}[caption={Własność obiektów pochodnych w kontenerze unikalnych wskaźników (C++14); kompletny program.}]
#include <cassert>
#include <memory>
#include <vector>

struct Operation {
    virtual ~Operation() = default;
    virtual int apply(int x) const = 0;
};
struct Add : Operation {
    explicit Add(int delta) : delta_(delta) {}
    int apply(int x) const override { return x + delta_; }
private:
    int delta_;
};
struct Multiply : Operation {
    explicit Multiply(int factor) : factor_(factor) {}
    int apply(int x) const override { return x * factor_; }
private:
    int factor_;
};

int main() {
    std::vector<std::unique_ptr<Operation>> operations;
    operations.push_back(std::make_unique<Add>(2));
    operations.push_back(std::make_unique<Multiply>(3));

    int result = 4;
    for (const auto& operation : operations)
        result = operation->apply(result);
    assert(result == 18); // (4 + 2) * 3
}
\end{lstlisting}

Warto pamiętać, że unikalny wskaźnik na typ pochodny można
przenieść do unikalnego wskaźnika na bazę tylko wtedy, gdy
spełnione są wymagania dotyczące konwersji wskaźników i deletera. Domyślny deleter w typowym
przypadku usuwa obiekt przez \texttt{Base*}, dlatego destruktor
polimorficznej bazy musi być poprawnie zaprojektowany.

\section{RTTI: \texttt{dynamic\_cast} oraz \texttt{typeid}}
\label{sec:cpp12-rtti}

RTTI (\emph{run-time type information}) pozwala w pewnych sytuacjach
sprawdzić typ dynamiczny i wykonać sprawdzoną konwersję w hierarchii
klas polimorficznych.

\subsection{\texttt{dynamic\_cast} — sprawdzony downcast}

\begin{lstlisting}[caption={Sprawdzona konwersja w dół hierarchii (C++11); kompletny program.}]
#include <cassert>
#include <typeinfo>

struct Animal { virtual ~Animal() = default; };
struct Dog : Animal { int bark() const { return 10; } };
struct Cat : Animal {};

int main() {
    Dog dog;
    Animal* animal = &dog;

    if (auto* ptr = dynamic_cast<Dog*>(animal)) {
        assert(ptr->bark() == 10);
    }
    assert(dynamic_cast<Cat*>(animal) == nullptr);

    Animal& reference = dog;
    assert(&dynamic_cast<Dog&>(reference) == &dog);
    try {
        (void)dynamic_cast<Cat&>(reference);
        assert(false); // nie powinno zostac wykonane
    } catch (const std::bad_cast&) {
        // Nieudany dynamic_cast do referencji rzuca bad_cast.
    }
}
\end{lstlisting}

W typowym \emph{downcast} źródłowy typ klasy musi być polimorficzny.
Nieudane rzutowanie do wskaźnika daje \texttt{nullptr}, a do referencji
zgłasza \texttt{std::bad\_cast}. Konwersja zwykłego wskaźnika
\texttt{Derived*} do dostępnej, jednoznacznej bazy \texttt{Base*}
nie potrzebuje RTTI ani \texttt{dynamic\_cast}.

\paragraph{\texttt{static\_cast} nie sprawdza typu dynamicznego.}

\begin{lstlisting}[caption={Niebezpieczny downcast; fragment celowo błędny.}]
struct Base { virtual ~Base() = default; };
struct Derived : Base { int only_derived = 7; };

Base base;
Base* ptr = &base;
// Kompilator moze zaakceptowac ponizsze rzutowanie,
// ale wskazywany obiekt nie jest obiektem Derived:
// Derived* wrong = static_cast<Derived*>(ptr);
// int x = wrong->only_derived; // niezdefiniowane zachowanie
\end{lstlisting}

\noindent Stosuj \texttt{static\_cast} przy downcast tylko wtedy,
gdy poprawny typ dynamiczny wynika z udowodnionego niezmiennika.
Jeżeli architektura wymaga bardzo wielu wywołań \texttt{dynamic\_cast},
rozważ lepszy interfejs wirtualny, visitor albo \texttt{std::variant}.

\subsection{\texttt{typeid} i typ dynamiczny}

\begin{lstlisting}[caption={Odczyt informacji o typie dynamicznym (C++11); kompletny program.}]
#include <cassert>
#include <typeinfo>

struct Base { virtual ~Base() = default; };
struct Derived : Base {};

int main() {
    Derived object;
    Base& base = object;
    assert(typeid(base) == typeid(Derived));
    assert(typeid(object) == typeid(Derived));
}
\end{lstlisting}

\noindent Dla wyrażenia typu klasy polimorficznej, będącego \emph{glvalue},
\texttt{typeid(expression)} zwraca informację o typie dynamicznym.
W innych sytuacjach zazwyczaj uwzględnia typ statyczny wyrażenia.
Nazwy zwracane przez \texttt{type\_info::name()} są zależne
od implementacji; nie należy używać ich jako stabilnego identyfikatora
w zapisywanych plikach lub w protokołach sieciowych.
Dla \texttt{typeid(*ptr)} z zerowym wskaźnikiem do typu polimorficznego
może zostać zgłoszony \texttt{std::bad\_typeid}.

\section{Wielodziedziczenie i wirtualne klasy bazowe}
\label{sec:cpp12-multiple}

C++ pozwala dziedziczyć po wielu klasach. Jest to przydatne
np. dla kilku niezależnych interfejsów, ale wymaga zrozumienia
niejednoznaczności oraz struktury podobiektów bazowych.

\begin{lstlisting}[caption={Dwie niezależne bazy-interfejsy (C++11); kompletny program.}]
#include <cassert>

struct Readable {
    virtual ~Readable() = default;
    virtual int read() const = 0;
};
struct Writable {
    virtual ~Writable() = default;
    virtual void write(int x) = 0;
};
struct Buffer : Readable, Writable {
    int value = 0;
    int read() const override { return value; }
    void write(int x) override { value = x; }
};

int main() {
    Buffer buffer;
    Writable& out = buffer;
    Readable& in = buffer;
    out.write(17);
    assert(in.read() == 17);
}
\end{lstlisting}

\subsection{Problem diamentu}

Załóżmy, że \texttt{Left} i \texttt{Right} dziedziczą po
\texttt{Base}, a \texttt{Diamond} po obu. Przy zwykłym
dziedziczeniu istnieją \emph{dwa różne} podobiekty \texttt{Base}.
Przy dziedziczeniu wirtualnym najbardziej pochodny obiekt zawiera
jedną współdzieloną wirtualną bazę \texttt{Base}.

\begin{lstlisting}[caption={Jedna wirtualna baza we wzorcu diamentu (C++11); kompletny program.}]
#include <cassert>

struct Root {
    explicit Root(int n) : value(n) {}
    int value;
};
struct Left : virtual Root {
    Left() : Root(1) {} // ignorowane jako inicjalizacja
};                          // bazy w obiekcie Diamond
struct Right : virtual Root {
    Right() : Root(2) {}
};
struct Diamond : Left, Right {
    Diamond() : Root(42), Left(), Right() {}
};

int main() {
    Diamond d;
    Root* from_left = static_cast<Left*>(&d);
    Root* from_right = static_cast<Right*>(&d);
    assert(from_left == from_right);
    assert(d.value == 42);
}
\end{lstlisting}

\noindent Najbardziej pochodna klasa \texttt{Diamond} odpowiada
za konstrukcję wirtualnej bazy \texttt{Root}. W implementacjach
wielodziedziczenie i wirtualne bazy mogą wymagać korekty wskaźników
oraz dodatkowych informacji o położeniu podobiektów.
Nie należy jednak utożsamiać modelu języka z konkretną reprezentacją
pamięci w danym ABI.

\section{Kowariantne typy zwracane}

W szczególnym przypadku nadpisująca funkcja wirtualna może zwracać
wskaźnik lub referencję do typu bardziej pochodnego, zamiast
wskaźnika lub referencji do klasy bazowej.

\begin{lstlisting}[caption={Kowariantny typ wyniku metody wirtualnej (C++11); kompletny program.}]
#include <cassert>

struct Base {
    virtual ~Base() = default;
    virtual Base* self() { return this; }
};
struct Derived : Base {
    Derived* self() override { return this; }
    int payload = 5;
};

int main() {
    Derived d;
    Base& b = d;
    Base* a = b.self(); // typ statyczny Base*
    Derived* x = d.self(); // typ statyczny Derived*
    assert(a == x);
    assert(x->payload == 5);
}
\end{lstlisting}

\noindent Kowariancja nie obejmuje dowolnych szablonów: na przykład
\texttt{std::unique\_ptr<Derived>} nie jest kowariantnym typem zwracanym
nadpisującym metodę zwracającą \texttt{std::unique\_ptr<Base>}.
Dlatego metody \texttt{clone()} zwracające \texttt{unique\_ptr}
w klasach pochodnych zwykle zachowują wspólny typ wyniku.

\section{Klasy abstrakcyjne jako kontrakty i wzorzec NVI}
\label{sec:cpp12-nvi}

Dobrą praktyką może być rozdzielenie publicznego, niewirtualnego
interfejsu kontrolującego niezmienniki od prywatnej implementacji
wirtualnej (\emph{non-virtual interface}, NVI).

\begin{lstlisting}[caption={NVI: publiczny kontrakt i prywatna metoda wirtualna (C++11); kompletny program.}]
#include <cassert>
#include <stdexcept>

class Encoder {
public:
    virtual ~Encoder() = default;

    int encode(int value) const {
        if (value < 0)
            throw std::invalid_argument("negative input");
        return do_encode(value);
    }

private:
    virtual int do_encode(int value) const = 0;
};

class DoubleEncoder final : public Encoder {
private:
    int do_encode(int value) const override {
        return 2 * value;
    }
};

int main() {
    DoubleEncoder encoder;
    Encoder& api = encoder;
    assert(api.encode(21) == 42);
}
\end{lstlisting}

\noindent Wersja publiczna może centralnie realizować walidację,
logowanie, kontrolę błędów i niezmienników. Prywatna funkcja
wirtualna w klasie bazowej \emph{może} zostać nadpisana przez klasę
pochodną, mimo że nie jest z niej bezpośrednio wywoływalna
przez zwykłą nazwę; w przykładzie nadpisanie jest prywatne.

\subsection{Polimorficzne kopiowanie przez \texttt{clone}}

Gdy chcemy uzyskać niezależną kopię obiektu o zachowanym typie
dynamicznym, możemy zastosować wirtualną fabrykę \texttt{clone()}.

\begin{lstlisting}[caption={Kopiowanie polimorficzne (C++14); kompletny program.}]
#include <cassert>
#include <memory>

struct Node {
    virtual ~Node() = default;
    virtual std::unique_ptr<Node> clone() const = 0;
    virtual int value() const = 0;
};
struct Number final : Node {
    explicit Number(int x) : x_(x) {}
    std::unique_ptr<Node> clone() const override {
        return std::make_unique<Number>(*this);
    }
    int value() const override { return x_; }
private:
    int x_;
};

int main() {
    std::unique_ptr<Node> a = std::make_unique<Number>(42);
    std::unique_ptr<Node> b = a->clone();
    assert(a.get() != b.get());
    assert(a->value() == b->value());
}
\end{lstlisting}

Takie podejście przywraca semantykę kopiowania do hierarchii,
której obiekty zwykle są trzymane przez wskaźniki.
Trzeba jednak jawnie zdefiniować, co oznacza \emph{głęboka kopia},
zwłaszcza jeśli obiekt współdzieli inne zasoby.

\section{Model implementacyjny: \emph{vptr}, \emph{vtable} i ABI}
\label{sec:cpp12-vtable}

Standard C++ określa \emph{zachowanie}, ale nie nakazuje użycia
żadnej konkretnej struktury wewnętrznej do realizacji funkcji
wirtualnych. Popularne ABI używają mniej więcej następującego modelu:
\begin{itemize}
\item obiekt polimorficzny przechowuje ukryty wskaźnik \emph{vptr};
\item wskaźnik ten prowadzi do tablicy \emph{vtable}, która zawiera
  informacje potrzebne do wywołań wirtualnych i często do RTTI;
\item wywołanie wirtualne może wymagać odczytania \texttt{vptr},
  odpowiedniego wpisu i wykonania wywołania pośredniego.
\end{itemize}

To \emph{opis typowej implementacji}, nie kontrakt języka.
W szczególności liczba i położenie wskaźników \texttt{vptr},
układ tablic oraz narzut pamięci są zależne od ABI i kompilatora.
Wielodziedziczenie może komplikować ten schemat, ponieważ niektóre
wywołania wymagają korekty wskaźnika \texttt{this}.

\begin{center}
\fbox{\begin{minipage}{0.91\textwidth}
\textbf{Schemat poglądowy (nie jest układem wymaganym przez standard)}
\medskip

\begin{tabular}{@{}p{0.43\textwidth}cp{0.36\textwidth}@{}}
\textbf{Obiekt \texttt{Derived}} & & \textbf{Tablica wirtualna} \\
\hline
\texttt{vptr} & $\longrightarrow$ & adres implementacji \texttt{draw()} \\
\texttt{Base::id} & & adres odpowiedniego destruktora \\
\texttt{Derived::payload} & & metadane zależne od ABI \\
\end{tabular}
\end{minipage}}
\end{center}

\subsection{\texttt{sizeof}, wyrównanie i brak gwarancji binarnego układu}

\begin{lstlisting}[caption={Empiryczny rozmiar klasy polimorficznej (C++11); kompletny program.}]
#include <iostream>

struct Plain { int x = 0; };
struct Poly {
    int x = 0;
    virtual ~Poly() = default;
};

int main() {
    std::cout << sizeof(Plain) << '\n';
    std::cout << sizeof(Poly) << '\n';
}
\end{lstlisting}

Nie wpisujemy w skrypcie stałych wartości tych rozmiarów:
na wielu platformach typ \texttt{Poly} będzie większy, ale
rozmiary wynikają z implementacji i obowiązujących reguł
wyrównania. Sama obecność \texttt{virtual} nie gwarantuje określonego
liczbowo narzutu ani ustalonej reprezentacji binarnej.
Nie należy serializować surowych bajtów obiektu polimorficznego
jako trwałego formatu danych.

\subsection{Koszty i devirtualizacja}

Wywołanie wirtualne może być droższe od zwykłego z powodu
pośrednictwa oraz potencjalnie trudniejszego \emph{inliningu}.
W praktyce jednak koszt zależy od rozkładu typów dynamicznych,
lokalności pamięci, optymalizacji międzyjednostkowej i architektury CPU.
Kompilator może zastosować \emph{devirtualizację}: jeśli potrafi
ustalić docelową implementację, wywołanie może stać się bezpośrednie,
a następnie zostać rozwinięte inline.

\begin{lstlisting}[caption={final może ułatwiać ustalenie celu wywołania; fragment.}]
struct Interface {
    virtual int eval(int) const = 0;
    virtual ~Interface() = default;
};
struct Square final : Interface {
    int eval(int x) const override { return x * x; }
};

int evaluate(const Square& x, int value) {
    return x.eval(value); // cel wywolania jest znany
}
\end{lstlisting}

\noindent Nie oznacza to, że \texttt{final} zawsze poprawia wydajność.
Porównania należy prowadzić na przykładach reprezentatywnych
i przy realistycznych opcjach optymalizacji. Wąskim gardłem może
być w istocie rozproszony układ obiektów w pamięci, a nie sam
koszt wywołania funkcji wirtualnej.

\section{Alternatywa: polimorfizm statyczny przez szablony i CRTP}
\label{sec:cpp12-static-poly}

Polimorfizm \emph{statyczny} wykorzystuje dedukcję typów i
instancjowanie szablonów, dzięki czemu kompilator zna konkretny typ
w miejscu użycia. Pozwala to generować kod wyspecjalizowany
bez konieczności utrzymywania wspólnej bazy wirtualnej.

\subsection{Szablon funkcji i concepts (C++20)}

\begin{lstlisting}[caption={Polimorfizm statyczny przez concept (C++20); kompletny program.}]
#include <cassert>
#include <concepts>

struct Circle {
    double radius;
    double area() const { return 3.141592653589793 * radius * radius; }
};
struct Rectangle {
    double a, b;
    double area() const { return a * b; }
};

template<class T>
concept HasArea = requires(const T& x) {
    { x.area() } -> std::convertible_to<double>;
};

template<HasArea T>
double area_of(const T& shape) {
    return shape.area();
}

int main() {
    assert(area_of(Rectangle{2, 3}) == 6.0);
    assert(area_of(Circle{2}) > 12.5);
}
\end{lstlisting}

Powyższy kod nie wymaga wspólnej klasy bazowej. Ograniczeniem jest to,
że w typowym kodzie szablonowym zestaw typów musi być znany przy
kompilacji miejsca instancjowania. Nie otrzymujemy automatycznie
heterogenicznego \texttt{std::vector<T>} dla wielu różnych typów.

\subsection{CRTP (\emph{Curiously Recurring Template Pattern})}

\begin{lstlisting}[caption={CRTP z metodą dostarczaną przez bazę (C++11); kompletny program.}]
#include <cassert>

template<class Derived>
struct Printable {
    int twice() const {
        const auto& self = static_cast<const Derived&>(*this);
        return 2 * self.value();
    }
};

struct Number : Printable<Number> {
    int data;
    explicit Number(int x) : data(x) {}
    int value() const { return data; }
};

int main() {
    Number n{21};
    assert(n.twice() == 42);
}
\end{lstlisting}

CRTP nie jest bezpośrednim zamiennikiem wszystkich interfejsów
wirtualnych: wyrażenie \texttt{Printable<A>*} nie jest uniwersalnym
wskaźnikiem do dowolnego \texttt{Printable<B>}. Może natomiast
umożliwiać \emph{static dispatch}, implementację domieszek
(\emph{mixins}) i wykrywanie błędów już przy instancjowaniu szablonów.

\section{Polimorfizm bez dziedziczenia: \texttt{std::variant}}
\label{sec:cpp12-variant}

Jeżeli lista możliwych typów jest zamknięta i znana przy kompilacji,
\texttt{std::variant} (C++17) pozwala przechowywać jedną z kilku
alternatyw \emph{przez wartość}. Zamiast wirtualnego wywołania
używa się zwykle \texttt{std::visit}.

\begin{lstlisting}[caption={Zamknięty zbiór typów i std::visit (C++17); kompletny program.}]
#include <cassert>
#include <variant>

struct Add { int x; };
struct Multiply { int x; };
using Instruction = std::variant<Add, Multiply>;

int execute(const Instruction& instruction, int input) {
    return std::visit([input](const auto& action) -> int {
        using T = std::decay_t<decltype(action)>;
        if constexpr (std::is_same_v<T, Add>)
            return input + action.x;
        else
            return input * action.x;
    }, instruction);
}

int main() {
    assert(execute(Add{2}, 4) == 6);
    assert(execute(Multiply{3}, 4) == 12);
}
\end{lstlisting}

\noindent To podejście ułatwia przechowywanie wartości bez osobnych
alokacji i ma inną charakterystykę wydajnościową. Ceną jest
konieczność jawnego wyliczenia alternatywnych typów w definicji
\texttt{variant} oraz aktualizowania odwiedzających, gdy pojawiają
się nowe alternatywy. W przykładzie użyto też
\texttt{std::decay\_t} oraz \texttt{std::is\_same\_v}
z nagłówka \texttt{<type\_traits>}.

\section{Polimorfizm dynamiczny i statyczny — porównanie}
\label{sec:cpp12-comparison}

\begin{center}
\small
\begin{tabularx}{\textwidth}{@{}p{0.23\textwidth}XXX@{}}
\toprule
\textbf{Cecha} & \textbf{\texttt{virtual}} & \textbf{Szablon/CRTP} & \textbf{\texttt{variant}} \\
\midrule
Zbiór implementacji & Otwarty w granicach hierarchii & Ustalany przy instancjowaniu & Zamknięty zestaw alternatyw \\
Wspólny uchwyt & \texttt{Base\&}, \texttt{Base*} & Zależny od typu & Jeden typ sumaryczny \\
Mechanizm wywołania & Zwykle dynamiczny & Zwykle statyczny & Wybór aktywnej alternatywy \\
Przechowywanie & Często wskaźniki do obiektów & Zazwyczaj konkretne obiekty & Wartość wewnątrz \texttt{variant} \\
Rozszerzanie o nowy typ & Zwykle bez zmiany wywołującego & Nowa instancjacja & Zmiana listy typów \\
Typowe koszty & Dyspozycja i lokalność danych & Rozmiar kodu po instancjacjach & Rozmiar największej alternatywy i odwiedzanie \\
\bottomrule
\end{tabularx}
\end{center}

Wybór zależy przede wszystkim od architektury programu:
\begin{itemize}
\item \textbf{Interfejsy pluginów i otwarte hierarchie}: zwykle
  polimorfizm dynamiczny, z uwzględnieniem stabilności ABI.
\item \textbf{Algorytmy generyczne na znanych typach}: zwykle
  szablony i concepts.
\item \textbf{Zamknięta rodzina wariantów wartości}: często
  \texttt{std::variant}.
\item \textbf{Wiele małych obiektów w intensywnej pętli}: przy wyborze
  należy zmierzyć także koszty alokacji, lokalność danych i wynik
  optymalizacji kompilatora.
\end{itemize}

\section{Reguła zera/pięciu w klasach polimorficznych}

Wprowadzenie destruktora wirtualnego powoduje, że automatyczne
generowanie operacji przenoszących nie działa tak samo jak w prostej
klasie agregującej pola. Dodatkowo publiczne kopiowanie samego
podobiektu bazowego może umożliwiać przypadkowe obcięcie obiektu.

W wielu hierarchiach interfejsowych rozsądne jest zabronienie
kopiowania i świadome zarządzanie obiektami przez smart pointery:

\begin{lstlisting}[caption={Niekopiowalna klasa bazowa interfejsu (C++11); fragment.}]
class Interface {
public:
    Interface() = default;
    virtual ~Interface() = default;
    Interface(const Interface&) = delete;
    Interface& operator=(const Interface&) = delete;
    Interface(Interface&&) = delete;
    Interface& operator=(Interface&&) = delete;

    virtual void run() = 0;
};
\end{lstlisting}

\noindent Nie jest to reguła uniwersalna: klasy reprezentujące wartości
mogą być kopiowalne, a wybrane hierarchie wspierają jawne
\texttt{clone()}. Istotne, aby polityka kopiowania, przenoszenia
i własności była świadomym elementem kontraktu.

\section{Typowe błędy i pytania kontrolne}

\begin{enumerate}
\item \textbf{Brak wirtualnego destruktora bazy}, mimo usuwania
  obiektów pochodnych przez wskaźnik do bazy.
\item \textbf{Brak \texttt{override}} i przypadkowe utworzenie innej metody
  przez pominięcie \texttt{const}, zmianę parametrów albo
  kwalifikatora \texttt{\&}/\texttt{\&\&}.
\item \textbf{Slicing} podczas przekazania przez wartość lub
  wstawiania do kontenera bazowych wartości.
\item \textbf{Bezpodstawny \texttt{static\_cast} w dół hierarchii}
  bez gwarancji rzeczywistego typu obiektu.
\item \textbf{Wywołanie funkcji czysto wirtualnej} podczas konstrukcji
  lub destrukcji na etapie, gdy implementacja pochodna nie jest aktywna.
\item \textbf{Mylenie implementacyjnej vtable z gwarancją standardu}.
\item \textbf{Zależność od domyślnych argumentów metod wirtualnych}
  różniących się między bazą a klasą pochodną.
\item \textbf{Nadmierna hierarchia klas}, gdy prostsza jest kompozycja,
  szablon lub \texttt{std::variant}.
\end{enumerate}

\noindent Pytania, które warto umieć rozstrzygnąć samodzielnie:
\begin{enumerate}
\item Czy \texttt{virtual} oznacza, że \emph{każde} wywołanie metody
  danej klasy musi być realizowane pośrednio?
\item Dlaczego \texttt{Base value = derived;} może dać inny wynik
  niż \texttt{Base\& ref = derived;}?
\item Co stanie się z funkcją wirtualną nadpisaną w \texttt{Derived},
  jeśli wywoła ją konstruktor \texttt{Base}?
\item W jakiej sytuacji \texttt{dynamic\_cast<Derived\&>} zgłasza
  \texttt{std::bad\_cast}?
\item Kiedy użyć \texttt{std::variant}, a kiedy interfejsu wirtualnego?
\end{enumerate}

\section{Wskazówki dotyczące diagnostyki i pomiarów}

Przy analizie klas polimorficznych przydatne są ostrzeżenia
\texttt{-Wall -Wextra -Wpedantic -Woverloaded-virtual} i, zależnie
od kompilatora, \texttt{-Wsuggest-override}. Błędy czasu życia
mogą być wykrywane przez AddressSanitizer i UBSan (w szczególności
odpowiednie kontrole vptr, jeśli są wspierane).

\begin{lstlisting}[language=bash,caption={Kompilacja przykładowego programu w GCC lub Clangu.}]
g++ -std=c++20 -O2 -Wall -Wextra -Wpedantic \
    -Woverloaded-virtual example.cpp -o example
./example

# Testy diagnostyczne (zalezne od platformy):
g++ -std=c++20 -O1 -g -fsanitize=address,undefined \
    -fno-omit-frame-pointer example.cpp -o example_san
./example_san
\end{lstlisting}

Należy porównywać z \emph{realnym} obciążeniem, a nie zakładać,
że jeden mechanizm jest zawsze szybszy od drugiego. Przy pomiarach
uwzględniamy optymalizację, lokalność danych, rozmiar próby,
rozrzut pomiarów i to, czy kompilator nie wyeliminował pracy.

\section{Podsumowanie rozdziału}

Kluczowe jest rozdzielenie trzech kwestii: \emph{jakie typy obiektów}
występują, \emph{kto odpowiada za ich czas życia} oraz
\emph{kiedy wybierana jest implementacja operacji}.
\texttt{virtual} umożliwia polimorfizm dynamiczny, lecz wymaga
rozważenia destruktorów i własności obiektów. Szablony i CRTP
oferują polimorfizm statyczny, a \texttt{std::variant} jest
wygodną alternatywą dla zamkniętej rodziny typów.
Praktyczna wydajność zależy nie tylko od kosztu wywołania,
ale również od reprezentacji i rozmieszczenia danych.

\section*{Literatura i dalsze źródła}
\addcontentsline{toc}{section}{Literatura do rozdziału 12}
\begin{itemize}
\item ISO C++ Standard: model obiektowy, klasy, dziedziczenie,
      funkcje wirtualne i RTTI (normatywne źródło reguł języka).
\item \url{https://en.cppreference.com/w/cpp/language/virtual}
      --- funkcje wirtualne i reguły nadpisywania.
\item \url{https://en.cppreference.com/w/cpp/language/dynamic_cast}
      --- sprawdzane konwersje w hierarchiach.
\item \url{https://en.cppreference.com/w/cpp/language/derived_class}
      --- dziedziczenie i wirtualne klasy bazowe.
\item Bjarne Stroustrup, \emph{The C++ Programming Language},
      wydanie czwarte, Addison-Wesley, 2013.
\end{itemize}
